%%%% Shortcuts
\newcommand{\num}[2]{\mbox{#1\,#2}}			% num with units

%%%% Symbols
\newcommand{\yes}{\ensuremath{\surd}\xspace}		% Tick mark
\newcommand{\no}{\ensuremath{\times}\xspace}		% Cross mark
\newcommand{\by}{\ensuremath{\times}\xspace}		% XXX x XXX
\newcommand{\bAND}{\ensuremath{\wedge}\xspace}		% Bool. /\
\newcommand{\bOR}{\ensuremath{\vee}\xspace}		% Bool. \/
\newcommand{\becomes}{\ensuremath{\rightarrow}\xspace}	% -->

%%%% Custom environments

% Centered tabular with single spacing
\newenvironment{ctabular}[1]
    {\par\begin{sspacing}\begin{center}\begin{tabular}{#1}}%
    {\end{tabular}\end{center}\end{sspacing}}


%%%% Our default level for display in TOC - subsubsections
\setcounter{tocdepth}{2}

% --- Glossary of key performance terms used in this thesis ---
% Throughput: number of completed requests per unit time (often denoted X).
% Service time S: time the system spends servicing a single request.
% Arrival rate \lambda: average number of arrivals per unit time.
% Utilisation U: fraction of service capacity used, typically U = X \times S.
% Response time R: total time from request arrival to completion (includes queueing and service delays).

